% RESEARCH_DRAFT_SHA256: 06c62e8bf5a7867511c1f8fac226ff26e91c45717efc5415aa31e7861f364119
% Layout and source-reference wrappers expanded for the web renderer.
\documentclass[11pt]{article}
\usepackage[margin=25mm]{geometry}
\usepackage{fontspec,unicode-math}
\setmainfont{Latin Modern Roman}
\setmathfont{Latin Modern Math}
\usepackage{amsmath,amssymb,mathtools,xurl,hyperref}
\hypersetup{colorlinks=true,urlcolor=blue,linkcolor=blue}
\setlength{\parindent}{0pt}
\setlength{\parskip}{5pt}
\setlength{\emergencystretch}{3em}
\newcommand{\sref}[2]{\hyperlink{source:#1:#2}{[S#2]}}
\newcommand{\cataloguescope}[1]{\paragraph{Source question.}#1}
\newcommand{\floor}[1]{\left\lfloor#1\right\rfloor}
\begin{document}

\section*{Zeta verification: bounded rational transcripts and a complete zero-dimensional certificate}
\noindent UnsolvedMath ID 20001265; AIM-COMPUTATION-0103. Research attempt dated 11 October 2026.
\paragraph{Result in scope.}This manuscript contains a source-specific attempt and complete proofs of the stated auxiliary calculations. It makes no new solution or novelty claim for the entire target.
\section{Original problem and preserved definition}
The following definition and source list are reproduced from the repository. Any earlier inspection dates and reported literature claims in this preserved text describe that source record, not additional verification by this attempt.
% BEGIN PRESERVED DEFINITION
\subsection{Definitions and mathematical statement}

\paragraph{Shared notation.}Standard number systems are denoted by $\mathbb N,\mathbb Z,\mathbb Q,\mathbb R,\mathbb C$; index ranges, computational models and specialized notation follow the cited problem source and the definitions below.


For a variety $V$ over $\mathbb F_q$, its zeta function is $Z(V,T)=\exp(\sum_{r\ge1}\#V(\mathbb F_{q^r})T^r/r)$. The decision task is whether a supplied function equals this zeta function. Membership in NP requires a finite certificate verifiable in time polynomial in the encoded input length.
\paragraph{Statement recorded in the supplied source.}
Question 15. Given a variety over a finite field, can you verify that a given function is its zeta function (i.e., is the question in NP)? (Lenstra, Edixhoven) \href{https://aimath.org/WWN/primesinp/primesinp.pdf}{[S2]}
\paragraph{Scope qualification.}
The source does not fix the encoding of the variety or candidate function. Its primary discussion mentions a fixed-genus case; the general target is retained without asserting that case settles it.
\paragraph{Source status.}
The supplied catalogue labels this record partially solved. Its source-status metadata was checked against the supplied snapshot on 2026-10-11. This import does not independently establish current open status or mathematical correctness.
\subsection{Short English statement}
Can a proposed zeta function of a finite-field variety be certified efficiently?
\subsection{Sources}
\begin{itemize}
\item[S1] Ulam AI and the UnsolvedMath Contributors, UnsolvedMath: A Curated Collection of Open Mathematics Problems, record 20001265 (AIM-COMPUTATION-0103). Catalogue curation and metadata: CC BY 4.0. The supplied source question is retained; definitions and the English summary are editorial additions.
\url{https://huggingface.co/datasets/ulamai/UnsolvedMath}.
\item[S2] AIM Workshop Problem Lists, original problem source linked by catalogue record AIM-COMPUTATION-0103
\url{https://aimath.org/WWN/primesinp/primesinp.pdf}.
\end{itemize}
% END PRESERVED DEFINITION
\subsection{Verbatim supplied catalogue question}
Question 15. Given a variety over a finite field, can you verify that a given function is its zeta function (i.e., is the question in NP)? (Lenstra, Edixhoven)

\section{A certificate route and its assumptions}
The original Question 15 asks about NP verification of a candidate
zeta function for a variety over a finite field. An NP claim requires
an encoding and certificate length polynomial in that encoding. The
argument below gives a rational-function reduction and a concrete
certificate for squarefree univariate zero-dimensional schemes. It
does not provide a certificate for arbitrary varieties.

\paragraph{How many point counts determine a bounded rational function.}
Suppose $Z_i=P_i/Q_i$ with $P_i(0)=Q_i(0)=1$,
$\deg P_i\le a$, $\deg Q_i\le b$, for $i=1,2$.
If their Taylor coefficients agree through degree $a+b$, then they
are equal. Indeed $P_1Q_2-P_2Q_1$ has degree at most $a+b$.
Multiplying $Z_1-Z_2$ by the series units $Q_1Q_2$ shows that
this polynomial has all coefficients through that degree zero,
so it is the zero polynomial. Conversely equality gives agreement.

For $Z(T)=\exp(\sum_{i\ge1}N_iT^i/i)=\sum_{k\ge0}z_kT^k$
over $\mathbb Q[[T]]$, differentiation yields
\[
kz_k=\sum_{i=1}^k N_i z_{k-i},\qquad z_0=1.
\]
The formal chain rule proves the identity after equating the coefficient
of $T^{k-1}$; division by $k$ recursively determines the coefficients.
Thus counts through $a+b$ determine $Z$ if the stated numerator
and denominator bounds are independently justified. They do not
prove those bounds or certify the counts themselves.

\paragraph{An explicit certificate in a nontrivial special case.}
Let $f\in\mathbb F_q[X]$ be monic, squarefree and of degree $D$.
Take $V$ to be its zero set as a zero-dimensional scheme. A certificate
is a list of distinct monic irreducible polynomials $f_j$ with
$f=\prod_j f_j$. If $d_j=\deg f_j$, the claimed answer is
\[
Z(V,T)=\prod_j(1-T^{d_j})^{-1}.
\]
Here is both its proof and a polynomial verification procedure.
For a monic polynomial $h$ of degree $d\ge1$, check
\[
h\mid X^{q^d}-X,
\quad
\gcd(h,X^{q^{d/r}}-X)=1
\quad\hbox{for each prime }r\mid d.
\]
These conditions characterize irreducibility. To prove this, the
roots of $X^{q^k}-X$ form the $q^k$-element field in an algebraic
closure: the polynomial has derivative $-1$ and degree $q^k$;
its roots are closed under field operations by the Frobenius identities.
A root of an irreducible degree-$e$ polynomial generates a field with
$q^e$ elements, since the quotient has a basis of $e$ powers of the
root. Its nonzero group has order $q^e-1$, so Lagrange's theorem
shows that its Frobenius orbit is finite. For a root $\alpha$ its
orbit under $x\mapsto x^q$ has length $e$
equal to its minimal-polynomial degree. The orbit product has coefficients
fixed by Frobenius, hence in $\mathbb F_q$, while every polynomial
over $\mathbb F_q$ vanishing at $\alpha$ vanishes on that orbit.
The degree equality follows. Such a root lies in the $q^k$-element
field exactly when its orbit length divides $k$.

The divisibility test therefore makes $h$ squarefree and forces every
irreducible factor degree $e$ to divide $d$. If $e<d$, some prime
$r$ divides $d/e$, so $e\mid d/r$ and that factor contradicts the
gcd test. All factors must have degree $d$, leaving just one because
$h$ itself has degree $d$. Conversely an irreducible $h$ of degree
$d$ satisfies all tests by the orbit characterization. Degree one
has no prime divisors to test and works by the same argument.

After verifying all factors, distinctness, and their product, a
degree-$d_j$ factor contributes $d_j$ points over $\mathbb F_{q^k}$
if $d_j\mid k$, and zero otherwise, again by the orbit description.
Its contribution to the exponent defining $Z$ is
$\sum_{m\ge1}d_jT^{md_j}/(md_j)=\sum_{m\ge1}T^{md_j}/m$.
Exponentiating gives $(1-T^{d_j})^{-1}$. This last formal identity
follows by differentiation and constant term one, or by the geometric
series. Multiplying the contributions proves the formula.

All tests have polynomial cost in $D$ and $\log q$ under an explicit
finite-field representation with polynomial-cost field arithmetic.
Repeated modular powering computes $X^{q^k}\bmod h$ using
$O(k\log q)$ polynomial multiplications by binary exponentiation;
$k\le D$. Euclidean gcd and dense polynomial products have polynomial
cost. The prime divisors of $d\le D$ can be found by trial division
through $D$, also polynomial in this parameter. The certificate lists
at most $D$ total factor degrees. A candidate rational function can
be compared with the displayed product by cross multiplication in
polynomial time in its input length and $D$.

\section{Where the general NP argument fails}
For an affine variety in $n$ coordinates, $N_k\le q^{kn}$, so
the integers $N_1,\ldots,N_B$ need at most
$O(nB^2\log q+B)$ bits. Short numerical answers are not automatically
short proofs. Listing every point to certify a count can already
require $q$ entries for the affine line over $\mathbb F_q$, exponential
in $\log q$. This only refutes that listing strategy; it does not
refute the existence of other NP certificates. Our squarefree-factor
certificate avoids enumeration in dimension zero, but no counterpart
is supplied for general varieties or for certifying general cohomological
degree bounds polynomial in the proposed input size.

\section{Audit}
The primary Question 15 was inspected. The audit distinguishes
determinacy, certificate soundness, certificate length and deterministic
construction of a certificate. It does not use a random factorization
algorithm to assert deterministic construction. The complete proofs
above concern the stated special input class, not the original universal
NP question. No factorization experiment or variety point count was run.

\section{Process and contributions}
Mathematics and model-assisted adversarial self-audit: OpenAI gpt-6-astra, ultra reasoning, as dispatched for this research batch; exact serving revision not exposed. Codex (GPT-6; exact serving revision not exposed) materially assisted with orchestration and format review. No human mathematical authorship, mathematical edits, or independent human proof checking is asserted.
Contribution prompt A was applied to the full definition and supplied source row, then prompt B was applied adversarially to the derived results. The item-specific audit above records the scope checks and rejected extrapolations. Local assembly and structural checks do not establish mathematical correctness. The source catalogue status is preserved in TOP\_PROBLEM and is separate from this attempt status.
\paragraph{Integration review.} Codex (GPT-6; exact serving revision not exposed) reviewed the written arguments and expanded layout and source-reference wrappers. Repository metadata and the production web renderer were checked using Python 3.13 and JavaScript with MathJax 3.2.2. This model-assisted review is not independent human proof verification; no PDF compilation or formal proof checking is asserted.

\section{Completion Estimate}
\noindent\textbf{COMPLETION: 0\%}
The proved material is a rederivation of known or auxiliary facts rather than a substantiated new advance on the original target. A complete certificate works for squarefree univariate zero-dimensional inputs; polynomial certificates for point counts and degree bounds of arbitrary varieties remain missing.
\section{References}
The complete original Sources subsection above is retained, including catalogue attribution and licence notices. The additional list records primary sources actually consulted for this attempt on 11 October 2026; uncited source links reproduced above were not all rechecked.
\begin{itemize}
\item American Institute of Mathematics, Future directions in algorithmic number theory, workshop 24--28 March 2003, report version 30 April 2003; original question, relevant lecture context and adjoining remarks inspected on 11 October 2026. \url{https://aimath.org/WWN/primesinp/primesinp.pdf}.
\end{itemize}
\section{Final status}
\textbf{Unresolved in this attempt.} A complete certificate works for squarefree univariate zero-dimensional inputs; polynomial certificates for point counts and degree bounds of arbitrary varieties remain missing.
\end{document}
